% Auto-generated from references/originals/enterprise-wechat-guide.pdf by scripts/bootstrap_sources.py.
% The extraction is intentionally conservative; revise wording and structure during editing.
\chapter{企业号关注流程及常见问题}
\sourceattribution{本章逆向整理自原 PDF《中国科学院大学企业号关注流程及常见问题解决方法》。原 PDF 未标注个人作者；截图和二维码请以归档原件为准。}

\section{关注流程}

\section{使用微信客户端扫描下方认证二维码或点击前往}

\section{根据页面提示输入账号密码进行第一步认证}

3.完善以下信息进行第二步认证，需填写微信绑定的手机号或邮箱（可在微信设置->账号与安全中查看）

注：此步骤可以更换手机号

\section{关注成功界面如下，如关注失败请点击查看常见问题}

\section{完成身份认证}

\section{常见问题及解决办法}

\section{确保自己在认证过程中填写的手机号是与当前使用微信绑定的手}

机号。进入微信，“设置 - 账号与安全 – 手机号”查看

\section{如何查看或设置微信号？}

\section{解绑}

已关注需要重新关注企业号进行认证，需先解绑微信再扫码关注进入个人信息页点击退出登录即可退出当前所在企业号(如果已经绑定微信，需要先解绑)。

\section{更换手机号}

（1）未关注需要更换手机号，可在认证到此页面时进行更换（2）已关注企业号需要更换手机号，可以先解绑企业微信，再扫码重新关注，按上步进行更换（3）出现手机号或邮箱不在企业通信录中是因为扫描的码不对，关闭页面重新扫描正确二维码进行关注，到此页面时可直接在通讯录中打开中国科学院大学企业号，查看是否关注成功

\section{认证失败的解决方法}

如果认证成功，识别下方二维码后，看不到应用，依然只有“身份认证小助手/企业小助手”一个应用时。可依次采用下面 3 种方法排查：

方法一

退出微信客户端，稍后重新进入查看。（腾讯服务器近期不太稳定，认证成功后的关注状态会有一定几率的延时）

方法二

返回“身份认证小助手/企业小助手”，点击“关注身份验证”卡片，输入认证过程中填写的手机号，进行认证。

若提示“该手机号不存在于企业通讯录中”，请尝试方法三。

方法三

若方法一、二无效，可尝试取消关注、扫描企业微信二维码重新关注，若设置正确，可直接关 注成功，无需身份认证。

取消关注的方式重新扫码关注

方法四

移动校园后台搜索学工号，进入个人信息编辑页，重新保存用户正确手机号，直接保存。

\section{认证失败解决方案（多次认证失败）}

用户认证时走到上图步骤，识别二维码进入，依然只能看到一个应用。

可能性一：用户认证时填写的手机号不是与用户当前微信绑定的手机号。

确认方式：1）请用户确认上图中“当前预留手机号”是不是和用户当前使用微信绑定的手机号2）进入微信【通讯录】-【新的朋友】，点击右上角【添加朋友】，输入上图手机号，点击搜索，将搜索到的微信头像和用户的微信头像比对，一致说明是，不一致说明手机号填写有误。

处理办法：

方法一：请用户输入和当前使用微信绑定的手机号进行认证

方法二：移动校园平台修改成正确的手机号，保存。用户即可关注成功（无需认证流程）

可能性二：微信服务器延时取消关注，扫描企业号二维码重新认证。

可能性三：

用户有 2 个手机号 x、y，分别绑定了不同的微信账号，已经用 x 手机号进行过身份认证，未解绑微信。用户又用 y 手机绑定的微信进行认证，认证时输入 y 的手机号，虽然页面提示关注成功，但因为已经用 x 手机号认证过，所以 y 手机绑定的微信依然看不到任何应用。

备注：这种情况下，y 手机进行认证时，在下面左图的页面里，系统会自动带出 X 的手机号，用户需要改成 y 手机号。不过，不管输入 x 还是 y 的手机号，y 手机上都是提示认证成功，但是 y 手机绑定的微信看不到任何应用。

A 进行身份认证，认证时输入 C 的手机号（C 手机号是与微信绑定的手机号，但不是 A认证时使用的微信）

如核实这种情况，按如下方法解决：

方法一：请用户输入正确的手机号，进行认证。

方法二：移动校园平台直接修改为正确的手机号，保存即可。

\section{账号密码试错解锁时间}

密码试错 6 次，账号会锁定，提示“错误次数已达最大上限，请稍后再试”解锁时间半个小时。

\section{取消企业微信关注的正确方法}

\section{外国手机号如何认证}

认证时，国际区号前补 2 个 0

\section{用户绑定的微信账号变更或与微信绑定的手机号变更了怎么办}

如果原微信账号还能登录：请用户进入微信个人中心页，解绑微信，用新手机号认证。如果原微信账号无法登录（例如手机号已注销、微信未解绑就换号了，导致旧的微信无法登录）：移动校园平台清除手机号和邮箱，保存，来解绑用户跟微信的绑定关系。

\section{企业微信如何解绑}

进入“身份验证小助手 - 身份认证”菜单的个人中心页面，点击“解绑微信”，再点击“解绑”解除跟微信的绑定关系

\section{多身份用户切换账号}

若多个学工号的证件号相同，用户可以进入个人中心切换身份，无需解绑微信，重新认证。

\section{报错提示}

手机号码已存在问题原因：（1）该手机号被人注册（2）该手机号是学生之前学号注册时绑定的，未解绑快速解决办法：由学校后台管理老师登录移动校园后台，搜索该手机号，确认信息后 将手机号从原账号下删除后保存，绑定到新账号下

\begin{editingnote}
原 PDF 包含二维码和多张微信、企业号操作截图。当前章节先保留文字逆向结果；排版时请对照 \path{references/originals/enterprise-wechat-guide.pdf} 补图，并核对 2026 级实际企业号认证流程。
\end{editingnote}
